\documentclass[11pt,twocolumn,titlepage]{article}
\usepackage{geometry}
\usepackage[review,endnotes]{legcite}
\usepackage{float}
\usepackage{indentfirst}
\usepackage{ifxetex}
\usepackage{fancyhdr}
\usepackage{color}
\usepackage{graphicx}
\usepackage{dblfloatfix}
\usepackage[bookmarksnumbered,pdfpagelabels,draft,implicit=false]{hyperref}

\geometry{
includeheadfoot,
hmargin=.75in,
top=0pt,
bottom=.4in,
headheight=.5in,
headsep=18pt,
foot=28pt
}

\ifxetex
    \usepackage{fontspec}
    \setmainfont[Ligatures=TeX]{Calibri}
\fi

\floatstyle{ruled}
\restylefloat{table}

\input citations

\title{A Five Part Plan for Patent Reform}
\author{Charles Duan}
\date{\today}

\sloppy
\hbadness=5000

\input pagestyle

\pagestyle{wp}

\newif\ifshortcompile
%\shortcompiletrue

\begin{document}

\onecolumn

\tracingoutput=1
\showboxdepth=1000
\showboxbreadth=1000

\thispagestyle{empty}
\hbox{}
\vfill
\hbox to \textwidth{%
    \hskip -1in
    \hskip 7pt
    \vbox to 0pt{%
        \vss
        \includegraphics[width=\paperwidth]{cover.pdf}%
        \vskip -7pt
        \vss
    }%
    \hss
}
\vfill
\hbox{}

\clearpage

\tracingoutput=0

\pagenumbering{roman}

\hbox{}
\vfill

\centerline{\includegraphics[width=.5\textwidth]{pkthinks.pdf}}
\vskip 0pt plus .3fill

\begingroup

\Large \parskip=.5\baselineskip \parindent=0pt
\leftskip=.5in \rightskip=\leftskip

The patent system is an important motivator for new technology. But in its
present state, the patent system has problems and loopholes that allow for abuse
and exploitation, harming rather than promoting innovation.

In this white paper, we look at how to fix those problems in patent law, by
identifying areas that are currently abused and that require reform. Those five
areas are:

\textbf{Accounting for all inventors.} The standards for patenting must account
for all types of inventors, large and small, and especially those who work
outside the patent system.

\textbf{Clarity of patents.} Patents are intended to disseminate knowledge about
new inventions and technologies. Thus, patent documents must be made clear and
understandable, so that they serve that function.

\textbf{Targeting the right parties.} Threats of complex patent litigation,
levied against consumers, small companies, and non-technology businesses, stifle
innovation without any corresponding benefit to inventors.

\textbf{Avoiding gamesmanship in litigation.} Licensing and enforcement of
patents should be about the merits of the patents, not about a party's
ability to run its opponents into the ground with litigation costs.

\textbf{Maintaining competition in the innovation economy.} Patent owners ought
to use their patents in ways consistent with long-standing principles favoring a
competitive marketplace.

\endgroup
\vfill

\hbox{}

\break

\tableofcontents
\clearpage
\twocolumn
\pagenumbering{arabic}

\parskip=.5\baselineskip plus .2\baselineskip minus .1\baselineskip

\section{Introduction}
\SectionNote
\AuthorityTag{Introduction}

Patents are simultaneously one of the most obscure and one of the most
contentious areas of the law. ``Even Rube Goldberg couldn't have invented
anything as confused and perverse as America's patent and intellectual property
protection system,'' said one scholar, comparing the patent system to the artist
``famous for crafting machines of ridiculous complexity.''\note{schrage} That
complexity has not stopped numerous efforts for reform, with members of Congress
introducing eleven patent-related bills in 2013.\note{see patentprogress-guide (listing the bills)}

One would think that this esotericity and contention would leave no room for agreement on patent policy. But there is one area of widespread agreement:
patents are one of the most abused areas of the law.
The term ``patent troll,'' coined by a lawyer
in the early 2000s, is
well known in the general lexicon to refer to an entity who produces no products
of its own but rather asserts patents, often patents on inventions that the
entity did not itself invent, against companies who do produce new technologies
and products.\note{iam-troll} Criticisms of patent trolls, and of the patent system that gives
rise to them, abound.

These accounts of abuses by patent trolls are neither hypothetical nor
academic. The harms to small businesses, innovators, and the economy are well
documented and empirically researched. One report documented that small
companies had stalled product development, delayed hiring new employees, or even
gone out of business as a result of patent assertion.\note{see chien-oti at 16 &
figure 3} Another widely-cited study estimated the direct costs of patent
assertion by non-practicing entities (that is, companies that exist solely to
assert patents)
at \$29 billion in 2011, with \$10.8 billion of that coming from
small- or medium-sized firms.\footnote{\sentence{see bessen-direct at table 4}.
The authors argue that much of that cost is deadweight loss, not contributing to
compensation for inventors.} Patent
litigation is big business, making it an attractive target to those wanting to
abuse it.

Although there are abuses, it is also necessary to account for the important
purposes and values
that underlie the patent system. By providing the reward of a patent,
patents give inventors
an incentive to research and develop new inventions, and also give those
inventors some lead time in bringing products to
market.\note{see, e.g., chakrabarty at 307 (``The patent laws promote [the
progress of science and the useful arts]
by offering inventors exclusive rights for a limited period as an
incentive for their inventiveness and research efforts.'');
bonitoboats at 146 (``The Patent Clause itself reflects
a balance between the need to encourage innovation and the avoidance of
monopolies which stifle competition\ldots.'')} Without patents, the theory goes, small inventors would never spend time
inventing, since large companies could sweep in, copy the inventions, and
undercut the small inventors' businesses. And certainly many inventions through
history, such as the light bulb, the airplane, and the photocopier, were
invented by small inventors and protected by patents.\note{see lemley at
722-723, 725-726, 734-735 (noting, in many cases, that many such inventions were
the result of incremental improvements made by multiple independent inventors or
teams of inventors)}

To find the right direction for policy reform of the patent system, it is
necessary first to reconcile these two competing threads, of the value patents
have created and the harm they are causing.

\subsection{A Generally Good System, But Often Abused}
\SectionNote

These two divergent views of the patent system, as destroyer of small
businesses and as engine of innovation, have led to wildly different and
apparently irreconcilable views of that patent system. On the one hand, some
have called the system fundamentally flawed and advocated for the abolition of
patents altogether.\note{see, e.g., boldrin at 12 (``Since there is no evidence
that intellectual monopoly achieves the desired purpose of increasing innovation
and creation, it has no benefits.\textellipsis This leads us to our final
conclusion: intellectual property is an unnecessary evil.'')} Others have
defended the need to protect inventors at all costs and opposed any changes that
might upset or potentially weaken patents.\note{see, e.g., mossoff-forbes
(``Weakening intellectual property laws due to negative policy rhetoric,
hyperbolic internet commentary, and extensive lobbying by firms who choose to
infringe patents because they don't want to pay the licenses offered to them by
patent licensing firms is irresponsible.'')}

But both of these views seem overly simplistic, and unfair to a complex system
such as patents. The question is whether there is an adequate middle ground,
where the benefits of the patent system can be maintained and even enhanced,
while the drawbacks can be removed.

At the most basic level, the answer to resolving these conflicting views on the
patent system seems to be this: there is a core of patents that is useful and
valuable to society. That core has been surrounded by layers and layers of
abusive practices, clever lawyering, and moneyed interests to turn the patent
system into something entirely different.

The question, then, is how to separate out that valuable core from the layers of
abuse. The simplest way to do so is to focus on the abuses that are currently
occurring, and to categorize those so that they may be attacked individually. It
is not difficult to find examples of those abuses, ranging from shotgun lawsuit
campaigns targeted against small retail and service businesses, to overbroad
patents stretched to cover basic ideas in technology, to monopolistic
arrangements intended to keep out competition. The remaining discussion, then,
will review these abusive practices within the patent system, and propose
reforms to curb those practices, within the framework of our five-part focus for
patent reform.

\begin{table*}
\vskip\baselineskip
\begin{center}
\begin{tabular}{|l|l|l|}
\hline
\textbf{Type} & \textbf{What covered} & \textbf{Exclusive rights} \\
\hline
Copyright & Creative works &
\rightskip=0pt plus 1fil\relax
Copying, distribution, public performance, etc. \\
Trademark & Product identifiers &
\rightskip=0pt plus 1fil\relax Use in commerce \\
Patent & Inventions &
\rightskip=0pt plus 1fil\relax Making, using, selling, etc. \\
\hline
\end{tabular}
\end{center}
\caption{Summary of major types of intellectual property.}
\bblabel[Table~\#]{iptypes}
\end{table*}

\subsection{Five Areas for Reform}

This white paper thus identifies five key areas in which the patent system is
ripe for reform. These areas were selected based on current patent reform
considerations, concerns raised by academics and the general public, and general
considerations of policy and consumer interests. The areas are:
\begin{itemize}
\item\textbf{Accounting for all inventors.} Patents are not the only stimulus
for invention and innovation. Numerous technology creators, such as computer
software startups, open innovation communities, and academia, have other
incentives to flourish,
discover, and advance science and technology. These
non-patent incentives should be celebrated, and not weakened by
overbroad protection of patents.
\item\textbf{Clarity of patents.} The grant of a patent is intended to disclose
knowledge in exchange for a temporary monopoly over a certain class of
inventions. This relies on clarity in two areas: (1) clarity in the description
of the invention, so that others may learn from it, and (2) clarity in the
description of the monopoly, so that others may know what is covered and what is
not. Where a patent uses imprecise language and fuzzy boundaries, that patent
may easily be abused and thus disserves the public.
\item\textbf{Targeting the right parties.} Traditionally, patents were the
domain of big technology companies, and the structure of patent acquisition and
enforcement grew around that domain. Today, however, patent lawsuits fall on the
doorsteps of small businesses and customers of technologies, parties who do not
expect to be part of the patent game and are ill equipped to play. This provides
an opportunity for abusers of the system to take advantage of unwary and
unsuspecting consumers.
\item\textbf{Avoiding gamesmanship in litigation.} Patent lawsuits are
complicated and expensive. Much of this complexity and expense is necessary and
expected, because the technologies are advanced, and the product markets
at stake are large. But clever litigants can exacerbate this complexity, hoping
to win cases not on the merits but rather on exhaustion and cost.
\item\textbf{Maintaining competition in the innovation economy.} Because patents
are a temporary monopoly, they necessarily and appropriately are an exception to
the general view that competition is preferable to monopoly. But the inclination
of any monopoly holder is to expand that monopoly beyond its anticipated reach.
Thus, patent owners have attempted to use patents to hold up technology,
block adoption of interconnection standards, and otherwise create undesirable
anticompetitive situations.
\end{itemize}
These five areas for reform are interrelated, and many proposed reforms will
deal with more than one area. But each of these five areas is independent, and
solving one will not automatically solve others. There is no
silver bullet to patent reform, and improving the patent system will require a
long-term, multifaceted effort. But that effort is worth the potential
innovation and creation that a well-crafted patent system will bring about.

\section{Background on the Patent System}
\AuthorityTag{Background}

Because the patent system is a fairly complex and unusual area of the law, some
background may be needed for those without much familiarity with patents. This
section will discuss the basic theory behind why patents are issued, and then
look at the mechanics of how patents are obtained and used.

\subsection{The Theory of Patents}
\SectionNote

``Intellectual property'' is the term applied to various legal systems in which
entities are given exclusive rights over non-tangible things. In the United
States, as summarized in \inline{this at iptypes},
there are generally three main types of intellectual property, which are
generally distinguishable by the type of non-tangible things to which the
exclusive rights are given. Copyrights cover creative works, such as books,
music, and artwork. Trademarks cover names that are attached to
products and services, such as brand names and logos. Patents cover inventions,
such as machines, processes, and systems. The focus here, of course, is patents.

The basic theory behind patents is that they provide an incentive for inventors
to invent. A person might come up with a great idea for a new machine for
shelling peanuts, say, but it might take a lot of time, money, research, and
development to turn that idea into a marketable product or service. But once the
shelling machine was on the market, then others could copy it and undercut the
price.\note{see denicolo at 366 (explaining that innovation is ``non-rival'' in
that ``once it is created, it can be shared at no cost'' by firms other than the
inventor)}

Without patents, the shelling machine inventor would probably do one of
two things: (1) not spend the time and money in developing the shelling
machine, choosing to do something else; or (2) operate the machine in a secret
factory, not allowing anyone to see it. Neither of these two options is
particularly desirable.\note{see fromer at 548 (``The theory is that this
stimulation occurs by rewarding inventors for taking two steps they likely would
not otherwise take: to invent in the first instance and to reveal information to
the public about these inventions.'')} The first one means that fewer new
inventions would be
created, and the second option means that other inventors will not be able to
improve on the invention.\note{see fromer at 548-549}

Thus, patents provide a way out of this dilemma. A patent offers the patent
owner the exclusive right to make, use, or sell the invention,\note{see 35usc at
S 271} as a \emph{quid pro quo} for
the inventor revealing to the world how the invention works in a written
document called a patent specification.\note{see, e.g., eldred at 225 (``The
issuance of a patent is appropriately regarded as a \emph{quid pro quo}---the
grant of a limited right for the inventor's disclosure and subsequent
contribution to the public domain.'')} Thus, inventors will have the incentives
to invest time and resources into inventing (solving the first problem), and the
public will have the benefit of the advancement in knowledge (solving the second
problem).\note{see, e.g., pennock at 19 (``{While one great object
was, by holding out a reasonable reward to inventors, and giving them an
exclusive right to their inventions for a limited period, to stimulate the
efforts of genius; the main object was `to promote the progress of science and
useful arts;' and this could be done best, by giving the public at large a right
to make, construct, use, and vend the thing invented, at as early a period as
possible; having a due regard to the rights of the inventor.}'')}

That last part is the key: patents are issued to advance the public's interest
in knowledge and access to new technologies. Patents should encourage inventors
to create and reveal new inventions for the benefit of everyone. This sentiment
is in fact embodied right in the Constitution of the United States, which
authorizes Congress to grant patents in order to ``promote the Progress of
Science and useful Arts.''\note{patent-clause} Even Thomas Jefferson, the first
patent examiner of
the U.S. Patent Office, saw ``the exclusive right to invention as given not of
natural right, but for the benefit of society.''\note{jefferson-letter}

\subsection{Obtaining and Using a Patent}
\SectionNote

To obtain a patent, an inventor submits a patent application to the U.S. Patent
and Trademark Office. The patent application includes two parts. The
\emph{specification}, which usually includes drawings and text, should describe
the invention in sufficient detail so that readers can recognize the inventor's
contribution and learn how to reconstruct it.\note{see 35usc at S 112/a} The
\emph{claims}, which are the
numbered paragraphs at the end of the patent document, are the legal
definition of the patent right.\note{see 35usc at S 112/b; osenga at 63-64
(``[Patent
claims] must serve the dual functions of putting boundaries on the patentee's
monopoly and providing sufficient notice of that monopoly to allow the public to
avoid infringing the patent.'')} Just as a title deed to a
plot of land specifies
the metes and bounds of the property, the claims specify the metes and bounds of
what infringes and what does not infringe the patent.\note{see mpep at
S 2173.05/a (``Applicants\ldots are required to make clear and precise the terms
that are used to define the invention whereby the metes and bounds of the
claimed invention can be ascertained.'')}

The Patent Office will then examine the application. This includes
reviewing the prior art: the examiner will search for similar
patents, technical papers, and other information that predates the application,
and try to find one that is sufficiently similar to the claims of the
application.\note{see 35usc at S 102/a (no patent should issue on an
invention that
``was patented, described in a printed publication, or in public use, on sale,
or otherwise available to the public'' before the relevant filing date); 35usc
at S 103 (no
patent should issue if ``the claimed invention as a whole would have been
obvious before the effective filing date'')}
If such prior art is found, then the examiner will reject the
application, and offer the applicant a chance to respond.\footnote{\sentence{see
mpep at S 706.02, 714}. The examiner may reject the application on other grounds
as well, such as indefiniteness of the claims or inadequacy of the disclosure.}
Often the examiner and the applicant will go through several rounds of exchanges
and revisions to the patent application, until the examiner is satisfied and
agrees to allow the application.\note{see mpep at S 1303} Then the Patent Office
issues a patent to the applicant.\note{see mpep at S 1309}

With patent in hand, the patent owner (called a ``patentee'') can now use the
patent, up until the patent term expires, to sue others who infringe the
patent.\note{see 35usc at S 281}
Another party infringes a patent if they make, sell, use, offer to sell, or
import something that comes within the claims of the patent.\note{35usc at S
271} Often the
patent
owner and the accused infringer disagree what the patent claims mean, so this is
resolved by the court hearing the lawsuit, in a procedure called ``claim
construction.''\note{see markman-fc at 978-979; osenga at 64-66; cand-patrules
at S 4 (claim construction proceedings)} Upon deciding
the meaning of the patent claims and other issues,
the court can decide whether the accused party infringes the patent. If so,
then the court can award money to the patent owner, order the infringer to stop
conducting the infringing activity, or both.\note{35usc at S 283-284}

\subsection{The Role of the Troll}
\SectionNote

Patent trolls are an arguably new phenomenon that have arisen in the last few
decades.\note{see, e.g., harkins at 410 (describing how a ``significant cottage
industry has grown up'' that ``ushered in the era of `patent trolls'\,''); but
see mossoff-myth} The term ``patent troll'' generally refers, to varying
degrees, to an entity that asserts patents without producing any products or
offering any services based on those patents. There is obviously a great deal of
uncertainty in what constitutes a patent troll, and authors have proposed
various alternate terms and definitions, such as ``non-practicing entity''
(NPE), ``patent assertion entity'' (PAE), and ``patent monetizing entity''
(PME), to name a few.\note{see ftc-2003 at chapter 2/(III)C2b(ii)B, at 31
(describing characteristics of NPEs); chien-arms at 328 (``Patent-assertion
entities are focused on the enforcement, rather than the active development or
commercialization of their patents.''); jeruss at 370 (defining patent
monetization entities as ones that ``described their main source of revenue as
patent litigation or licensing'')}

The rhetoric around patent reform often focuses on the harms of patent trolls
and suggests a need to target those entities and their practices. And, indeed,
there are good reasons to believe that when an entity does not actually put its
patents into use, that entity is substantially more likely to engage in
undesirable and abusive behaviors. For example, entities that do not practice
their patents lack the incentive to truly educate the public about their
inventions,\note{see duffy at 1395 (``Holding all else equal, a practiced patent
discloses more, teaches more, and contributes more to the sum total of social
knowledge than does a mere paper patent.'')} but rather they may prefer to wait
for others to come up with those inventions independently and then spring up
demanding a license fee. Also,
patent trolls lack business relationships with manufacturers and other industry
players, and are thus free to engage in abusive business practices
without fear of repercussions.\note{see lemley-forest at 2165 (``It seems likely
that practicing entities have in the past been more concerned than trolls about
such reputational matters.'')}

But patent trolls are symptomatic of broader problems with the patent system
overall. Because the system permits abusive practices, it is not surprising
that a business model has grown to take advantage of and profit from those
abuses. But those abuses can continue to exist and to be used whether or not
patent trolls are present---and, indeed, some evidence suggests that even
operating, product-producing companies engage in the same sorts of abusive
practices ordinarily ascribed to patent trolls.\note{see lemley-forest at 2170
(``When it comes to patent enforcement, it is hard to conclude that trolls cost
society more than practicing entities.'')}

Thus, efforts to reform the patent system ought not to focus on defining and
opposing patent trolls as a business model. Rather, efforts should be targeted
at rooting out the loopholes, problems, and abuses of the system overall. When
patent trolls are given nothing to take advantage of in the patent system, they
will naturally wither away, and we will be left with a stronger system that
promotes innovation and progress.

\section{Accounting for All Inventors}
\SectionNote
\AuthorityTag{Accounting for All Inventors}
 
The traditional rationale behind the issuance of patents is utilitarian: by
granting limited monopolies on inventions, the government can incentivize
individuals and firms to spend resources on inventing. The limited monopoly,
namely the right to exclude others from practicing the invention for a period
of time, both gives a financial reward to those who invent and grants an
opportunity for them to commercialize their inventions without interference
from competitors.\note{see, e.g., fromer at 547 & note 31 (citing sources)}

The empirical evidence for this rationale is mixed. The area in which the
patent incentive most clearly succeeds is in the pharmaceutical
industry,\note{see eisenberg at 350 (``Biopharmaceutical research is often held
out as a shining example of
the success of the patent system in motivating private investment in
R\&D.'')}
but some have suggested that this is primarily the case because of artificially
imposed
regulatory requirements that necessitate the exclusive lead time offered by
patents.\note{see eisenberg at 346 (noting that patent advocates cite
drug regulation ``as a large part of the cost of drug development that can only
be recovered if firms are allowed to charge patent-protected premium prices for
new products.''); but see eisenberg at 359 (arguing that FDA regulation of new
drugs in fact acts as a ``pseudo-patent'' complementing patent exclusivity)}
In other fields of technology, surveys and other evidence have
suggested that the patent incentive is insubstantial or vastly
ignored.\note{see, e.g., seltzer at 929 (``Patents do not provide a useful
incentive to innovate in the software industry, I contend, because the patent
promise ill-suits the engineering and development practices and business
strategies of software production.''); lemley at 736-738; checkpoint-brief at 3
(``Software patents are not necessary to spur innovation among the Amici. Our
engineers do not innovate because they hope to get patents.'')}

In any event, it is unnecessary to consider in too much depth the empirical
evidence for the patent incentive, because it is certainly not the only
incentive for innovation. The fast-paced startup community, open source
software, the prestige of publication, and prizes and rewards for new
discoveries, among other things,
are all alternative
incentives for innovation beyond patents, as explained below.\note{see kewanee
at 484 (``Certainly the patent policy of encouraging invention is not disturbed
by the existence of another form of incentive to invention.'')}

\subsection{Computer Software Startups}
\SectionNote

The widespread use of the Internet and computer
technologies has led to a proliferation in small software startups. The
market of mobile device applications, for example, is globally valued at over
\$53 billion\note{pappas} and accounts for approximately 466,000 jobs created
since 2007.\note{mandel at 13}
Many of the major Internet companies today have their roots in such a small
business.\note{see menell at 389 (``Many of the notable innovations---from
Napster's peer-to-peer system to the Google search engine, YouTube, and
Facebook---were hatched and initially developed in the Web 2.0 equivalent of the
Silicon Valley garage: a dorm room.''); seltzer at 973 (noting that the software
industry features ``{lower capital costs\ldots; lower pure research costs and a
focus on implementation; and lower uncertainty of development}'')}

Small software technology startups often do not rely on patents for protecting
their innovations. With patent applications costing about \$25,000 to file
and prosecute,\note{lemley-vc at 138 note 3 (citing that cost as ``not
unreasonable based on my experience''); see also ipwatchdog-cost (citing costs
of \$5,000 to \$15,000 for filing the application alone)}
they are well beyond the financial reach of these small startups
that may not even have \$25,000 to run their business in the first
year.\footnote{The startup accelerator Y Combinator, for example, invests
\$20,000 for a three-person startup. \sentence{see wired-ycombinator}.}
Yet the
number of such startups is proof that there is sufficient incentive to innovate
in that arena even without patents.

Instead, different incentives predominate. The first mover advantage is
particularly valuable in the fast-paced world of computer software.\note{seltzer
at 982-983} Network
effects hold strongly with software, as adopters of a company's
technology are less likely to move to competitive offerings due to built-up
social networks, cost of converting data, and familiarity with user
interfaces.\note{seltzer at 983-984}
Additionally, software companies are prone to failure for numerous reasons, so many software engineers are comfortable
with frequent ``pivots'' to entirely new ideas.\note{see ries at 149 (describing a pivot as ``a structured course correction designed to test a new fundamental hypothesis about the product, strategy, and engine of growth'')}
The mere experience of starting a
software company, say many such engineers, is incentive enough to innovate even
in the face of daunting odds, as it is easy and nearly cost-free to abandon one
idea and move on to the next.

\subsection{Open Innovation Communities}
\SectionNote

Open innovation communities are collectives of individuals and entities who
openly share their innovations, making those innovations available
to others for use, adaptation, and improvement.\footnote{The term ``open
innovation community'' is from \clause{dpl at 2 note 1}.} Although the open
source software community is perhaps the best known of these, open innovation
communities may be found in many other areas of technology besides software
development, including electronic hardware manufacturing, 3D printing, biology,
and environmental engineering.\footnote{\sentence{see dpl at 2 & note 1}.}

To strong proponents of the patent incentive theory, open source software and
related models of innovation present a quandary: how can innovation occur in a
world where products are given away for free and competitors are allowed---even
encouraged---to copy? Yet widespread use and constant improvement of open source
software suggests that those incentives must still be present. Some of the most
widely used software programs today, including the GNU/Linux operating system,
the Apache HTTP server, and the Firefox web browser, were developed by the open
source community.

Indeed, scholars have documented those alternative incentives that have
contributed to the growth of the open source software and other communities.
Reputation benefits play a significant role: as one seminal work put it in
describing two popular software projects, ``by properly rewarding the egos of
many other hackers, a strong developer/coordinator can use the Internet to
capture the benefits of having lots of co-developers.''\note{cathedral-bazaar at
43 (describing Linux and fetchmail); see also dpl at 17} Companies like IBM and
Red Hat invest in open source development to accrue returns such as consulting
services.\note{see benkler at 46; brodkin-red-hat} And basic ideals of sharing
and disseminating knowledge motivate others.\note{see, e.g., stallman at 129
(``My work on free software is motivated by an idealistic goal: spreading
freedom and cooperation. I want to encourage free software to spread, replacing
proprietary software that forbids cooperation, and thus make our society
better.'')} Thus, a variety of incentives, entirely apart from the patent
incentive, can spur innovation within open innovation communities.

\subsection{Alternative Rewards for Invention}
\SectionNote

Patents are one type of reward for innovation, but there are many others.
Prizes for innovation have been suggested as a solution to the economic
inefficiency of patent monopolies. ``The alternative of awarding prizes would be
more efficient and more equitable,'' writes one prominent
economist.\note{stiglitz; see also kremer (proposing an auction process by
which the government buys out most patents, offering the buyout price as a
prize)} And  numerous prizes are granted to incentivize societal progress:
Alfred Nobel, for example, bequeathed his fortune to establish prizes awarded to
``those who, during the preceding year, shall have conferred the greatest
benefit to mankind.''\note{nobel-will}

There are many examples of rewards for innovation, beyond the straightforward
prize. Academics are rewarded for their ideas and discoveries by having their
papers accepted in journals. Foundations run competitions for the first person
to solve an unsolved problem, to encourage inventors to develop creative
solutions. Governments provide tax incentives for research and development. And
researchers with new ideas can apply for grants, both government and
privately-funded, to pursue those ideas.\footnote{For a comprehensive overview
of rewards for innovation, see generally \clause{hemel}.}

An example of the last of these types of rewards incentivizing innovation may be
found in the development of the Internet. Many Internet technologies were built
under federal grants,\footnote{Much of early Internet technology research was
funded by the Advanced Research Projects Agency (ARPA, now DARPA). \sentence{see
rfc2235 at 1-6}.}
which incidentally restricted acquisition of
patents on the technology.\note{See 35usc at S 200-212}
Further Internet development was supervised
by standards organizations like the W3C, which expressly disallows patenting of
technology adopted into standards.\footnote{More specifically, if a participant
in a W3C standard acquires a patent essential to the standard, the participant
must generally grant a royalty-free license to the public for that patent.
\sentence{w3c-patent at S 3.1, 5}.}
The reward of a federal grant or
incorporation into an Internet standard was sufficient to bring about incredible
technological development without the use of patents.

Certainly none of these alternate incentives can entirely supplant the patent
system, as each person is motivated by different incentives. What is important,
though, is that no single type of incentive is the sole or predominant engine of
new innovation and technology.

\subsection{Conflicts Between the Incentives}
\SectionNote

While, in an ideal world, all of these
incentives would complement each other to maximize invention, in reality these
incentives conflict and sometimes work at opposing purposes. The focus here is
particularly on how the patent incentive can conflict with other incentives to
innovate, and how reforms to the patent system can reduce these conflicts.

For example, the first mover advantage is an incentive for startup
entrepreneurs, and that incentive is undercut when second movers or even
non-movers acquire and assert patents. Thus, there is substantial anecdotal
evidence of innovative software startups dropping products or closing shop
altogether in
the face of patent threats.\note{see, e.g., chien-oti at 16-17 & figure 3; tal
at 6:20-6:37 (``[Interviewer]: Did it put your business in danger? [Startup
founder] Jeff Kelling: It did, and they knew that. The settlement they wanted to
get was just enough to put us in danger, but not to close us, and I'll stop
there.'')}

The interference between patents and open source software is also well known.
One study from 2004 has shown that the Linux kernel, a popular and widely-used
open source program, potentially infringes 283 patents.\note{osrm-linux at 1}
Indeed, when a bundle
of 882 patents were proposed to be sold in 2011, the Department of Justice
intervened out of concern that the patents would ``jeopardize the ability of
open-source software, such as Linux, to continue to innovate and
compete.''\note{doj-cptn}
Similarly, the authors of the GNU General Public License are of the view that
patents ``obstruct free software development,'' which led to the inclusion of a
mandatory patent license in the most recent version of the GPL.\note{whygpl;
gplv3 at S 11}

There is also a conflict between the patent incentive and incentives of
alternate rewards. Patents can interfere with the ability of scholars to conduct
research, thereby diminishing the ability of academia to pursue innovative
ideas. However, this example is instructive, because patent law already
provides an accommodation for this conflict: an exception for experimental use.
That doctrine, which exempts from patent infringement non-commercial
experimental uses of patented inventions,\footnote{\bblabel{expuse}\sentence{see
whittemore at 1121 (``[I]t could never have been the intention of the
legislature to punish a man, who constructed such a machine merely for
philosophical experiments.''); rosuvastatin at 527 (``However, patenting does
not deprive the public of the right to experiment with and improve upon the
patented subject matter.''); see also cls-enbanc at 1322-1325 (Newman, J.,
concurring and dissenting) (discussing experimental use doctrine)}.}
accommodates those who invent to
obtain public or academic recognition. Additionally instructive, however, is
the degree to which this experimental use doctrine has been narrowed over
time,\note{see eisenberg-exp at 1019-1020 (``For the most part, the courts have
held that the experimental use defense does not apply to the facts of the
particular cases before them.'')}
reflecting an unfortunate shift in the relative valuation
between the patent incentive and rewards-based research incentives.

\subsection{Finding the Right Balance}
\SectionNote

One task for patent reform, then, is to consider
adjustments to the patent system that better accommodate these alternate
incentives for innovation. The goal of such adjustments is to better encourage
these inventors incentivized by factors other than patents, and to ensure that
patents do not stand in the way of those inventors.

Consider the following ideas for recalibrating the patent system in view of
these alternate incentives and alternate inventors.

For one thing, the standards for granting patents should be reconsidered. Many
aspects of the current patent system reflect an assumption that patents are the
primary driver of innovation, and modern developments challenge that assumption.
This question of patent quality is discussed in \inline{this at quality}.

Along similar lines, patent law includes an old doctrine called ``experimental
use,'' which protects those who engage in experimentation from the risk of
patent infringement.\note{see this[discussion] at expuse and accompanying text}
However, over the years courts have sharply narrowed the
availability of the experimental use exception, exactly at a time when more and
more experimentation is occurring, on the part of consumers, home inventors, and
other small parties. Expansion of this doctrine could provide protection to that
grassroots innovation that could otherwise be threatened by patents.

Additionally, some have considered the possibility of using specially-configured
patent licenses to promote open innovation rather than interfere with it. The
idea is
that a properly crafted license agreement attached to a patent could encourage
others to share knowledge rather than hold it closed behind patents, in the same
way that copyright licenses like the GPL and Creative Commons licenses encourage
others to share creative works.

One such proposal, the Defensive Patent License (DPL), institutes a system by
which an entity may agree not to assert its patents against other DPL members,
and in exchange the entity receives an automatic, free license to all other
patents under the DPL.\note{dpl at 38-39} Thus, the DPL ``blends the general
strategy of defensive patenting with the [open innovation community] values of
openness and freedom'' to ``provide an interoperable, technologically neutral,
reliable, and legally binding commitment to defense.''\note{dpl at 38} Proposals
like these take the patent system, which is built around the singular patent
incentive, and turn it around to incentivize other values shared by different
kinds of inventors.\footnote{For another example of such an alternate patent
license, see \clause{twitter-license}.}

\subsection{Improving Patent Quality}
\SectionNote
\bblabel[Section \#]{quality}

A common complaint about the patent system is the number of ``low-quality''
patents that are issued.\note{see, e.g., ftc-2003 at chapter 4/(II)A3, at 8
(``Several participants voiced concern about too great an issuance of obvious
patents.'')} There are frequent reports of patents on old or obvious ideas,
particularly in the field of software.\note{see, e.g., eff-podcasting} Many
people believe that these poor quality patents are the root of the problems that
the patent system faces today. They thus call for better examination of patents
through increased funding to the Patent Office, better training of examiners,
and longer time for examination.\note{see generally adams-iam (collecting
numerous views on patent quality)}

Improving patent quality is a key component of accounting for all inventors and
innovators. The process of invention is incremental,\footnote{``The world goes
ahead because each of us builds on the work of our predecessors. `A dwarf
standing on the shoulders of a giant can see farther than the giant
himself.'\,'' \sentence{sony at 477 note 28 (quoting \clause{chafee at 511});
see also lemley at 714; newton (``If I have seen further it is by standing on
the sholders \sic! of Giants.'')}.} so inventors depend on a pool of knowledge
not encumbered by patents, on which their inventions may be based. Part of this
pool comes from unpatentable abstract ideas, laws of nature, and physical
phenomena, which the Supreme Court has described as ``part of the storehouse of
knowledge of all men\ldots free to all men and reserved exclusively to
none.''\note{bilski at 3225 (quoting \clause{funkbros at 130}) (omission in
original)} Another part of this pool comes from old technologies in the prior
art,\note{see 35usc at S 102} as well as the knowledge that would be obvious to
a
``person having ordinary skill in the art.''\note{35usc at S 103} If a
patent erroneously issues on a technology within that public pool, or if courts
misread the boundaries of the pool too narrowly, then future innovators and
future innovation are hindered.\note{see aptea at 152-153 (an obvious patent
``withdraws what already is known into the field of its monopoly and diminishes
the resources available to skillful men'')}

The need to account for all inventors extends to many
areas of the patent quality problem. Patent examination is one such area. The
Patent Office traditionally searched for prior art primarily in the library of
past patents,\note{see mpep at S 902.03(e); p2p at 135} but today, where so much
innovation occurs outside of the patent system, the Patent Office is starting to
seek out external sources of information to discover prior art.\footnote{There
have been several formal efforts toward externalizing, or ``crowdsourcing,'' the
prior art review effort. \sentence{see p2p-pr; 35usc at S 122/e1 (preissuance
submissions by third parties in pending patent applications); wh-crowdsourcing
(``[T]he USPTO is announcing a new initiative focused on expanding ways for
companies, experts, and the general public to help patent examiners, holders,
and applicants find relevant `prior art'\ldots.'')}}

The basic standard for patent quality is another area where accounting for all
inventors must play a role. The Supreme Court recognized this in the case
\inline{ksr}.\note{*} For many years, the courts had used a test for what would
be obvious to a person of ordinary skill in the art, namely that some
``teaching, suggestion, or motivation'' to combine two different ideas must
exist before the combination could be deemed obvious.\note{ksr at 1734} The
Supreme Court rejected that test, based on its reasoning that the ``diversity of
inventive pursuits and of modern technology counsels against limiting the
analysis'' to that test---that is, because the test failed to account for all
types of inventors and technology.\note{ksr at 1742} Nevertheless, the
requirement for some sort of motivation to combine even persists after
\inline{ksr}.\note{see, e.g., kinetic at 1368-1369 (reversing a finding of
obviousness for a patent due in
part to a lack of expert testimony on motivation to combine)}

That conflict between the ``motivation to combine'' test and actual inventors is
placed sharply in focus when considering open-source software developers, for
example. Non-proprietary software developers and other innovation communities
value interoperability and combinability of software. Thus, the legal assumption
that new combinations are uncommon and often worthy of patents conflicts with
the experiences of those software developers, for whom new combinations are
routine and expected. Uncorrected, this mismatch means that patents would likely
stifle rather than encourage the tinkering and exploration that drives many
innovators today.

A critical step in improving patent quality, then, is accounting for the ways in
which all types of inventors work and innovate. The Patent Office and the courts
need to know how a ``person having ordinary skill in the art'' approaches
problem solving and invention. Without sufficient contact with actual innovators
of all types, these decisionmaking bodies will not be able to craft patent
policy that promotes innovation for all.

So it is important for these decisionmakers to reach out to inventing
communities, even those that do not file for patents, and it is important for
those communities to reach out to the Patent Office and other
decisionmakers.\note{see quality-comments at 4-6} Indeed, the White House and
Patent Office have initiated several programs to obtain input on patent
policy,\note{see, e.g., crowdsourcing-roundtable} and these initiatives should
be continued and broadened to reach the whole innovating community.

\section{Clarity of Patents}
\SectionNote
\AuthorityTag{Clarity of Patents}

Among the most common complaints about the patent system is that patents are
impossible to read. Patent documents are filled with complex language that
often, to the lay reader or technical engineer, appears confusing or even
misleading.\footnote{This section is drawn in part from
\clause{nautilus-brief}.}

It is no surprise, then, that those with technical but not legal skill do not
find patents to be a useful source of technical information. One survey of
researchers in the nanotechnology field found found numerous complaints on the
usefulness of patents as technical disclosures.\note{ouellette} Among the
surveyed researchers, 36\% never read patents, many of them complaining about
``the style in which patents are written---patents were called `vague,' `legal
jargon,' `incomprehensible,' and lacking `technical detail.'\,''\note{ouellette
at 571} Those who had read patents had similar complaints: ``To a scientist,''
an academic chemist wrote, ``the patent literature looks like an invention of
lawyers for the benefit of other patent lawyers.''\note{ouellette at 576} Only
38\% of surveyed researchers believed that one could reproduce an invention from
the patent specification---a clear indication that disclosure and dissemination
of knowledge, the cornerstone goals of the patent system, are not being
served.\note{ouellette at 576}

Those outside of academic circles have expressed the most concern about the
difficulty of understanding patents and particularly patent claims. One author,
in describing how to read patent claims, laments that they are ``a dense form
of legal English,'' further explaining that the ``drafting of patent claims is a
black art'' because claims are ``full of magic words.''\note{weir}
``Patents are so vague,'' one attorney said, that ``[i]f someone claims
infringement, the only way to resolve it is through
litigating.''\note{steeves} Where litigation is ``the only way'' to assess
the meaning
of a patent, that patent has failed to serve its public notice function.

This lack of clarity in patents is particularly unfortunate given the intention
that patents are meant to serve, namely as disclosures to the world of new
inventions and discoveries. This intention is embodied in the Constitutional
provision authorizing the grant of patents, for ``promoting the Progress of
Science and the useful Arts.''\note{patent-clause}  That intention of disclosure
further underlies many of the traditional utilitarian justifications for
patents.

There is, of course, no reason that patents could not be written to be clear.
The same sorts of inventions have been described in trade publications and
academic journals for centuries, where clarity of explanation is prized.
Instead, unclarity in patents arises out of the legal environment in which
those patents are born.

For example, the claims of patents are often the most criticized parts of
patents when it comes to lack of the clarity. The claims of a patent are
supposed to be written ``particularly pointing out and distinctly claiming'' the
invention,\note{35usc at S 112/b} so that they inform the public of what is
covered by the
patent and what is not.\note{see mcclain at 424 (holding that a patent claim
secures ``all to which [the patentee] is entitled'' while ``appris[ing] the
public of what is still open to them''); unitedcarbon at 236 (finding that
ambiguous patent claims create a ``zone of uncertainty'' that will ``discourage
invention''); ge at 369 (``The limits of a patent must be known
for\textellipsis the assurance that the subject of the patent will be dedicated
ultimately to the public.'')} However, a series of cases has whittled away at
these principles, permitting even highly vague and ambiguous claims to stand
valid so long as they are not ``insolubly ambiguous,'' even if ``reasonable
persons will disagree'' over the meaning.\note{exxon at 1375; accord datamize at
1437} This test fails to serve that crucial public notice function: as one
patent attorney put it, ``if reasonable people can disagree about the definition
of the claim terms, how does this notify the public of the patentee's right to
exclude?''\note{hartung at 844; see also laser at 27 (``If the Federal Circuit
does not correct this trend soon, a competitor's ability to accurately determine
the metes and bounds of current patents might deteriorate further.''); lemley at
745 (``[T]he Federal Circuit has permitted a number
of vague general disclosures that don't actually communicate very much to
anyone, and patent lawyers often have incentives to write such vague
disclosures.'')} The situation is equally dismal when it comes to searching for
patents of relevance to a particular company.\note{mulligan-lee at 297-298}
The permissibility of such unclear patents thus impedes operating businesses
from safely operating, while providing no societal benefit in return.

The rules of patent interpretation sometimes encourage this unclarity. As one
example, due to various court cases that have read poorly worded patent
specifications as limiting the scope of patents, it is common practice to
include long lists of alternatives. Consider the following definition of
``electronic network'' from a patent:
\begin{quote}
     Electronic Network---an electronic communication medium across which
     sellers
     and buyers may communicate, especially when communicating through the
     owner's main site. Representative electronic networks include the Internet,
     intranets, the public switched telephone network (``PSTN''), wireless voice
     and data networks, and television networks, such as satellite, broadcast,
     cable television, and two-way interactive cable. Electronic networks
     further include hybrid systems, such as those in which sellers communicate
     to buyers via one medium, such as cable television, and buyers communicate
     to sellers via another medium, such as the Internet. Electronic networks
     additionally include aggregated electronic networks, such as when buyers
     communicate to sellers via multiple media, such the Internet, the
     telephone, and cable television.\note{627patent at column 4, lines 26-40}
\end{quote}
This definition provides no useful technical information, and is in ways
contradictory to the ordinary understanding of ``electronic network'' (one
usually does not categorize the telephone network as such), so it thus serves
only to attempt to broaden the scope of the patent, at the expense of others
being able to understand it easily.\looseness=1

More importantly, though, is the paradoxical situation that the most valuable
patent is the most ambiguous patent. A patent that is prone to multiple
interpretations will cost a targeted defendant more to analyze, thus making it
more likely that the defendant will settle a case for a nuisance amount, and
also making it more likely that any litigation over that patent will be costly
and protracted. Thus, one economist traced how patents with ``fuzzy boundaries''
have created ``a business opportunity based on acquiring patents that can be
read to cover existing technologies and asserting those
patents.''\note{bessen-regulation at 34}. Similarly, stakeholders reported to
the
Federal Trade Commission that the patent system ``generally creates `an
incentive to be as vague and ambiguous as you can with your claims' and to
`defer clarity at all costs.'\,''\note{ftcreport at 85; see also ppi at 3
    (citing ``vague or expansive terms'' in patents as one of ``three cross
    winds'' creating a flood of patent litigation); segal (``But as long as the
    [patent] system exists, [noted patent assertor Erich] Spangenberg is going
    to exploit its ambiguities and pokiness for all it's worth.''); bessen-book
    at 164 (``The evidence suggests~.~.~.~that the deterioration of the notice
    function might be the central factor fueling the growth in patent
litigation'')}

The patent examination process can furthermore be exploited to exacerbate patent
unclarity. Often the patentability of an invention turns on the particular
wording of the claims that define that invention, and the examiner's
interpretation of that wording.\note{see phillips at 1317
(``[T]he prosecution history provides evidence of how the PTO and the inventor
understood the patent.'' (citing \clause{lemelson at 1206}))} Although the
examiner's interpretation ought to
be set forth in writing in the record
of proceedings (the ``file history'' or ``file wrapper'') for any
patent,\note{see cfr1-2 (``All business with the Patent and
Trademark Office should be transacted in writing.''); cfr1-133 at /b (``In every
instance where reconsideration is requested in view of an interview with an
examiner, a complete written statement of the reasons presented at the interview
as warranting favorable action must be filed by the applicant.'')}  in
practice much of that information is never recorded.\note{see phillips at 1317
    (observing that the file wrapper is ``less useful for claim construction
    purposes'' because it ``often lacks the clarity of the specification'')} Patent
    applicants, for example, can make oral presentations to
examiners (called ``interviews'' in patent practice), in which the applicants
explain aspects of their invention and the claims.\note{see cfr1-133; mpep at S 713} Much of this explanation is
never recorded in the file history, leaving the now-owner of the patent free to
assert any other interpretation of the patent that benefits the
owner.\note{eff-comments at 12 (``Summaries [of interviews] are often too vague
    and brief to provide meaningful notice to the public. The lack of a full
    record means that applicants can make arguments in interviews without facing
    the estoppel issues that would accompany a written filing.''); see also
    ipfrontline (``Another possible reason to argue the merits of a patent
    application orally rather than in writing is to avoid placing statements in
    the record that can be used against a patentee years later\ldots. Although
    the USPTO rules require patent applicants to provide a written record of an
    interview to be placed into the application file, statements in such written
    records are much more limited.''); inventivestep-interview (patent attorney
    noting that, for some arguments presented in an interview, ``I may omit it
    from my next response, thereby keeping it out of the record''); cf. mpep at
    S 713.04 (``A complete written statement as to the substance of any
    face-to-face, video conference, electronic mail or telephone interview with
    regard to the merits of an application must be made of record in the
    application, whether or not an agreement with the examiner was reached at
    the interview.'')}

Improving the clarity of the patent system is a multifaceted task, requiring
support from many different areas.

\subsection{Judicial Interpretations of Patent Language}
\SectionNote

Courts interpret patents during a
phase of trial litigation called ``claim construction,''\note{see, e.g.,
cand-rules at SS {4{-}1} to {-7}} and as such have the
final say in how patents are interpreted.\note{see markman at 391
    (holding that interpretation of claim terms is ``an issue for the judge, not
    the jury''); cybor at 1456 (``[A]s a purely legal question, we review claim
    construction de novo on appeal including any allegedly fact-based questions
relating to claim construction.'')} Thus, it is important that good law be
issued from these courts, and especially the Federal Circuit, the court of
appeals that oversees and makes many of the rules for claim construction.
Among other things, courts must carefully
scrutinize patents without falling victim to obfuscating techniques used to
make patented inventions appear more complex, broader, or more innovative than
they actually are.\footnote{For an example of this situation, see
\clause{wildtangent-brief}.}

\subsection{Patent Examination Procedures}
\SectionNote

Since the Patent Office is the first body to
interpret the text of patent applications, it is critical that its reasoning be
made clear and available to others wishing to understand the content of
patents.

Part of this effort is simply to build a better record of the Patent Office's
interpretations as it examines applications. Currently, examiners write Office
Actions that identify reasons for rejecting applications, most often over prior
art references. However, those Office Actions do not generally explain in
detail how the examiner read the patent application or interpreted particular
words in the claims.\note{see phillips at 1317 (noting that an Office Action
``often lacks the clarity of the specification and thus is less useful for claim
construction purposes'')}
Such information would benefit the public in later reading
the patent, and also assist the patent applicant in understanding the basis for
rejection. Furthermore, communications between the examiner and the applicant
should be made of record.\note{see glossary-comments at 4}

Additionally, the Patent Office should push for applicants to draft clearer
patents. One of the requirements for granting a patent is drafted claims
``particularly pointing out and distinctly claiming the subject matter which the
applicant regards as his invention.''\note{35usc at S 112/b}
Patent examiners could more rigorously
enforce this requirement, asking applicants to choose clearer, simpler language
that more distinctly defines the invention.\footnote{The Patent Office is
currently looking at a ``pilot program aimed at encouraging the use of clearer
language within patent claims.'' \sentence{wh-crowdsourcing}.}

\subsection{Limiting the Volume Game}
\SectionNote
\bblabel[Section \#]{volume}

One strategy for manufacturing unintelligibility in
patents is increasing volume: applying for patents with hundreds of claims, and
using continuation practice to acquire dozens of patents on the same
invention.\note{see moore at 81}
This volume game contributes nothing to the public store of knowledge---an
invention is secured through a few claims just as it is secured through
hundreds---and the purpose of acquiring this thicket of patent claims is to
increase the burden on those trying to assess the nature of the patents and
operate businesses in clearance of them.\note{see shapiro at 126 (``In short,
with multiple overlapping patents\ldots we have a volatile mix of two powerful
types of transaction costs that can burden innovation\ldots.''); moore at 82-83
(noting particularly that such continuation practice with respect to
pharmaceutical patents ``serves no useful social purpose'')}

Limiting this practice of unnecessary inflation of patent portfolios would
contribute much to streamlining the patent system, simplifying the process for
applicants, examiners, and third parties. This reform could be implemented
through Patent Office rules, through legislation, and/or by judicial
rule.\footnote{The USPTO attempted to institute certain limiting rules on
continuations, but that effort was rejected by the Federal Circuit, \clause{see
tafas}, and the USPTO declined to pursue the effort further, instead voluntarily
rescinding the proposed rules, \clause{see tafas-settlement}.}

\subsection{Improved Technological Tools}
\SectionNote

Although patent language is complex and
technical, it is also highly structured, potentially making it amenable to
automated analysis such as natural language processing.\note{see, e.g.,
sheremetyeva} Such systems could
simplify the interpretation of patents and greatly reduce the cost of having to
hire experienced counsel to review patents. Additionally, automated systems
would introduce a level of objectivity into patent interpretation, and
potentially stave off incentives for gamesmanship in obscure patent
language.\looseness=1


\section{Targeting the Right Parties}
\SectionNote
\AuthorityTag{Targeting the Right Parties}
 
A truly astonishing aspect of the patent litigation seen today is its volume:
dozens of companies sued at a time, with hundreds or thousands more the targets
of demand letters. Many of the targets are not the manufacturers of supposedly
infringing products, but rather the users. Thus, small businesses offering
wireless Internet access have been threatened for purchasing WiFi
routers,\note{innovatio}
iPhone
application developers have been threatened for using Apple's in-app purchase
mechanism,\note{lodsys}
and innumerable individuals and companies have been targeted simply
for having a scanner in the office.\note{scanner-troll}
Today's patent wars are against end users
and consumers.

\subsection{Nineteenth Century Patent Trolls}
\SectionNote
\bblabel[Section \#]{farm-patents}

Historically, this simply was not the case: patent battles, while still
prevalent, were between large companies, not end users of technologies. The
patent wars over sewing machines, airplanes, and automobiles, for example,
involved only small numbers of patents and parties to lawsuits.\note{see
chien-sw at 336-337}

In one historical case, end users of technology were targeted. During the
late nineteenth century,
numerous patents were granted over agrarian tools, which led to ``scores of
individual farmers who were sued for infringement based on farming tools they
had bought, rather than invented.''\note{chien-sw at 330} Contemporaries at the
time complained that ``the country is so large and the number of articles under
patent so great that there are not one-tenth of our farmers who know whether
their implements are patented or not.''\note{hayter at 65} As a result, patent
owners and attorneys with ``\,`bully' and `wily' methods often convinced the
inexperienced and `innocent' farmer that the patent covered the exact tool or
implement in question and collected a royalty fee.''\note{hayter at 66}
Ultimately, this led to a change in the law that made such patents much more
difficult to obtain,\note{see chien-sw at 348} and associations of farmers also
fought and won cases invalidating some of the most egregious patents, on basic
farm tools like swinging gates and drivewells.\note{see hayter at 73-74, 76-77}

Today we are seeing similar assertion of patents against end users because
technology has become ubiquitous. Just as swinging gates and drivewells were
necessary and widely used tools for an agrarian society in the 1800s, cell
phones, computers, and software are necessary and widely used tools for every
business today.

As a result, parties who ordinarily would never be involved in the patent
system are being dragged into it, due to their ordinary use of software
technology.
The Government Accountability Office found, in a study, that 39\% of patent
lawsuits involving software patents were against non-technology firms:
\begin{quote}
One representative from a retail company noted that historically, all of the
patent infringement lawsuits brought against the company used to be related to
products they sold. However, as of mid-2012, the representative said that half
of the lawsuits against the company were related to e-commerce software that the
company uses for its shopping website---such as software that allows customers
to locate their stores on the website---and were brought by [patent monetization
entities]. Representatives of retail and pharmaceutical companies told us they
also defend lawsuits brought by PMEs related to features on their
websites---typically software that outside vendors provide to them, rather than
something they developed. Additionally, city public transit agencies have been
sued for allegedly infringing patents by using software for real-time public
transit arrival notifications, according to a few stakeholders we
interviewed.\note{gao-report at 23}
\end{quote}

These entities are attractive targets for lawsuits for a number of reasons, none
of which is beneficial to the overall patent system. For one thing, end users
tend to be unfamiliar with the patent system, so they are poorly equipped to
defend themselves, compared to the technology companies that have dealt with
such lawsuits traditionally.\note{see love at 1614 (``Also, compared to
customers, manufacturers have a relative advantage litigating patent suits
because they generally have greater knowledge of the industry, the prior art,
and the patented invention's value.'')} Additionally, once end users have
incorporated the manufacturer's product into their products or businesses,
juries tend to unintentionally inflate the value of the patent due to the
unrelated features of the products or services.\note{see love at 1633-1634
(``[T]he larger the accused device, the harder it is for jurors to distinguish
between value attributable to the patented invention and value attributable to
other features and components.''); chao at 122 (discussing the
``overcompensation problem'' that is ``particularly true when the patent
involved only covers a single component of a multiple component product'')}
These factors give abusive patent asserters the upper hand, without any
corresponding benefit to the public or to innovation.

Finding the right solution for protecting end users is difficult. The ordinary
case---a customer buys a product off the shelf and is sued for patent
infringement---is straightforward and sympathetic, but there is a multitude of
business arrangements between technology manufacturers and users. End users may
simply use a product, may modify the product, or may incorporate it into another
device. Manufacturers may make a single line of products, or multiple lines, or
custom products for each customer. These and other factors demand a nuanced
solution to the problem of abuses against end users of technology. Several
proposed solutions are discussed below.

\subsection{Redirecting Patent Suits to Manufacturers}
\SectionNote

The most direct reform would be to redirect suits against end users toward
manufacturer suits in appropriate situations. Some, for example, have suggested
providing immunity to patent lawsuits for end users of technology who simply
purchase products off the shelf.\note{long (citing the Electronic Frontier
Foundation's statement on legislative solutions for patent reform); rackspace
(``Congressman Farenthold offered a simple, workable solution to solve a big
part of the problem: immunize end users of commodity-type products from patent
litigation.'')} The patent owner would remain able to pursue claims against the
manufacturer. This would satisfy fair compensation
for the patent owner while ensuring that the patent dispute remains between
parties experienced and expecting to be involved in patent lawsuits.

The difficulty with direct immunization is that there are situations
in which compensation from the end user is appropriate.\note{see, e.g.,
patentlyo-surplus (suggesting an economic model in which theoretically
compensation from an end user would be appropriate); but see patentlyo-surplus
(noting that, in reality, the aforementioned analysis, ``even if mathematically
correct, is largely irrelevant in the current climate'')}
For example, when the
end user is a technology company who commissions a product from a small
developer, it may be appropriate for the technology company to defend a patent
infringement lawsuit. As a result, such an immunity would need to
be carefully tailored with appropriate exceptions to ensure that the burden of
defense is placed on the appropriate party in all situations.

One attempt to overcome this problem would be to permit lawsuits against end
users but allow the manufacturer to be brought into the lawsuit as appropriate.
For example, one scholar suggests an arrangement in which the end user could
pull in the manufacturer into the lawsuit.\footnote{\sentence{see chao at
132-134}. Chao adopts theory of contributory liability from torts, under which,
if two people cause an injury and one of those two is sued, the sued person can
bring in the other person. ``If contribution were applied to patent law,
manufacturers of multi-component products accused of patent infringement would
be able to demand that any supplier of an infringing component share in any
potential liability.'' \sentence{chao at 133}.} This relieves some pressure on
the end user, since the manufacturer is now involved in the lawsuit, but it
would not solve the problem of lawsuits against customers who lack the resources
to even defend themselves in the first place.

\subsection{Customer Suit Exception}
\SectionNote

A second option deals with the timing of lawsuits against end users, rather than
the substance of the lawsuit itself. The principle, known as the ``customer suit
exception,'' is that if both the manufacturer and the end user are involved in
separate patent lawsuits, then in appropriate situations the lawsuit against the
end user will be ``stayed,'' or held off, until the manufacturer suit
completes.

The customer suit exception already exists in the law.\note{katz-v-lear at 1464
(explaining that, under the customer suit exception, ``litigation against or
brought by the manufacturer of infringing goods takes precedence over a suit by
the patent owner against customers of the manufacturer'')} Courts justify the
exception based on two reasons. First, if a decision is reached on the
manufacturer suit, it will affect all the customer suits, but if a decision is
reached on one customer suit, it will not necessarily affect other customers, so
the manufacturer suit will save on overall litigation costs.\note{see love at
1616-1617 (``In its first few decades of existence, courts applied the customer
suit exception relatively liberally, justifying its application on efficiency
grounds\ldots.'')} Second, courts recognized that, even in a suit against a
customer, ``in reality, the manufacturer is the true defendant in the
customer suit.''\note{see codex at 737-738; love at 1617}

However, over time, courts have cut back on the customer suit
exception.\note{see love at 1617-1618} Thus, there are several legislative
proposals on the table to revive and expand the exception.\note{see goodlatte at
section 5, S 296; leahy at section 4, S 299A; farenthold at section 4, S 300}
Doing so would help to balance the rights of patent owners and end users of
products.\note{love at 1635}

\iffalse
\subsection{Calibrating Damages}
\SectionNote
\bblabel[Section \#]{damages}

The attractiveness of bringing suit against end users
arises partly out of the idiosyncratic methods of calculating damages for
patent infringement. Often courts attempt to compute a ``reasonable royalty''
intended to reflect a percentage of revenues that the patent owner would have
successfully negotiated as a license fee. However, this amount is often grossly
overestimated: a traditional rule of thumb, for example, was that 25\% of
[revenues/profits] were due as a license fee. Even though the 25\% rule has
since been disapproved, courts still routinely award ``reasonable'' royalties
amounting to large chunks of a product's value.

Excessively high royalty rates are a particular problem for end users. If the
royalty rate is too high in relation to the value the end user derives from the
product, then the end user would be dissuaded from purchasing the product,
creating an unproductive chilling effect on the adoption of technology and
ultimately harming the consuming public.

The problem is magnified because the value an end user derives from a product is
hard to quantify. A consumer purchases a product for a variety of reasons and
put it to a variety of uses, and often the patented aspect of the product is
less important or unimportant to the consumer. Consider, for example, a coffee
shop that installs a WiFi router for its customers' use, and is sued for patent
infringement over a certain aspect of wireless network technology. How does one
attribute a portion of coffee sales to the patent, which only indirectly adds
value to the WiFi router, and which even more indirectly influences the
purchases of coffee?

Accordingly, as discussed below,\note{see this at reasonable-royalty} if damages
for patent infringement were more carefully calibrated
to the actual value of the patent vis à vis the business of the targeted
company, suits against end users would likely diminish greatly because patent
owners would stand to win little from such end users.
\fi

\subsection{Abolishing Software Patents Is Insufficient}
\SectionNote

A common thought is that, since
many of the problems described above (and throughout this paper) arise in the
field of software patents, the simple solution is to eliminate patentability of
software.\note{see, e.g., lee (``Reiterating that "pure" software can't be
patented wouldn't just be good law — it would also save the nation billions of
dollars in litigation costs.''); burton at 88-90 (finding, in an informal survey
of 49 software programmers, that 29 of them (59\%) thought that the United
States should abolish software patents)}

There are some practical issues in actually eliminating software patents
(devising a definition of software patents that cannot be easily circumvented by
clever lawyers is a problem\note{see chien-sw at 354 (`` One of the biggest
challenges to `abolishing software patents' is the question of what exactly is a
`software patent'?''); but see chien-sw at 356 (suggesting that ``a working
definition, rather than a perfect definition, may be what is really needed'')}),
but the more fundamental problem is that reforms specifically targeted toward
software patents are temporary fixes for the problems of today.
Although software patents are
substantially problematic now,\note{see, e.g., gao-report at 22 (``Specifically,
about 84 percent of [patent monetization entity] lawsuits from 2007 to 2011
involved software-related patents\ldots.''); bessen-switch}
this is likely because software is one of the
most widely used technologies now.\note{see, e.g., brynjolfsson at 10 (``The
difference between being a winner and being a lagging firm in IT-intensive
industries is very large and growing. Using technology effectively matters more
now than ever before.'')}

Analogously, as discussed above, farm technology was
widely used in the nineteenth century,
and patents on farm technology were hotly contested.\note{see this at
farm-patents}
Patents on those farm tools were effectively abolished.\note{see chien-sw at
348}
But that fix to the patent system did not prevent the software patent
problems faced today---it ultimately was a Band-Aid rather than a cure.

The same would be true of eliminating software patents.
The fundamental issue is that the technologies of
tomorrow are unknown,\footnote{As an example, as the 3D printing industry grows,
so does interest in patents on 3D printing. \sentence{see, e.g., wee-3dp
(describing the ``patent land grab of 3-D intellectual property'')}.}
so targeting patent reform to one specific field of
technology means that the same problems will only arise again in a different
technological sector.


\section{Avoiding Gamesmanship in Litigation}
\SectionNote
\AuthorityTag{Avoiding Gamesmanship in Litigation}

Patents are a privilege given by the government, and as such they must be
enforced
through the channels of government. In the United States, this means that
enforcement of patents must be through the federal district courts.\note{see
35usc at S 281} Thus, an
essential part of the patent system must be an efficient litigation mechanism
for enforcing patent rights.

However, recent history has shown a patent litigation system that is used less
and less for enforcement of legitimate rights, and more and more for abusive
tactics intended to enrich a few savvy litigants while impoverishing the
technology-consuming public.\note{see, e.g., ppi at 5 (``By filing plausible
claims raising sufficient risk of loss for defendants, [PAEs] force companies to
choose between expending millions of dollars to litigate the claims or settle
for less than the cost of the litigation.\ldots [M]ost patent assertions are
meant to drive settlements.'')}
The ease of engaging in these abusive tactics, and
the profitability of doing so, has given rise to an entire industry of patent
litigants, known variously as non-practicing entities, patent assertion
entities, or (derisively) patent trolls.\note{see ftcreport at 50 note 2, !60
(distinguishing these terms)}.

Low-quality patents are often blamed for the patent troll problem.\note{see,
e.g., gao-report at 31-32 (``Some of the representatives from operating
companies also said that [patent monetizing entities] are often more willing to
bring lawsuits based on a broad interpretation of their patents'
claims\ldots.''); mershon-trolls (``Critics say patent trolls amass vague or
low-quality patents and use them to sue or extract licensing fees from
companies, putting a strain on businesses and the larger economy.'')}
And patent assertion entities often assert low-quality patents.\footnote{As one
study found, nonpracticing entities owned 63.5\% of the most-litigated patents.
\sentence{alw at 708}. The study further found that NPEs won only 8\% of their
cases on the merits, while operating companies won 40\% of their cases.
\sentence{alw at 708}.}
But there are numerous
tactics that are entirely unrelated to the quality of the patents being
asserted, tactics that drive up costs of litigation and force unjustified
settlements or court awards. It is these tactics that are the subject of this
section.

Abusive practices can be found through the entire patent enforcement process,
from pre-litigation communications to post-trial appeals. The following are a
few prominent types of abuses, and ways to address them.

\subsection{Demand Letters}
\SectionNote

Prior to any lawsuit being filed, a patent owner can send a letter, arguing that
the recipient of the letter infringes a patent and demanding royalty
payments.\note{see, e.g., swang (``Demand letters are letters patent owners send
to unsuspecting businesses or individuals alleging patent infringement and
threatening a lawsuit.'')}
In itself, the process of sending a demand letter is perfectly ordinary, and
even commendable insofar as it is preferable to negotiate an arrangement without
having to incur the costs of court procedures.

Recently, however, a number of unscrupulous patent owners have found an easy way
to abuse demand letters. They send letters to small companies, who lack the
resources or experience to evaluate the merits of the arguments presented. They
then are free to make outrageous or inflated claims, demanding payments without
even explaining what products infringe the patents, hoping that the letter
recipients are too scared or cash-strapped to fight back.

Some of these abusive demand letters even go so far as to mislead or deceive the
recipients. Many threaten immediate litigation, when the sender actually has no
intention of bringing a lawsuit. Some accuse recipients of infringement even
when the sender has clearly not done any investigation to determine if the
recipient actually infringes the patents asserted.\note{see generally
hec-testimony at 3-5 (discussing these various abusive demand letter practices)}
These misleading and
deceptive practices have raised eyebrows, and triggered investigative efforts at
both the federal and state levels.\note{see vtne-ag at 3-5}

Today little is known about the world of patent demand letters, due in no small
part to the efforts of the senders of those letters, who use nondisclosure
agreements and shell companies to hide their activities.\note{see hec-testimony
at 6; ftc-study at 61{,}353 (``While workshop panelists and commenters identified
potential harms and efficiencies of [patent assertion entity] activity, they
noted a lack of empirical data in this area, and recommended that the Federal
Trade Commission use its authority\ldots to collect information on PAE
acquisition, litigation, and licensing practices.'')}
Thus, tackling the problem of abusive demand letters must begin with learning
about the scope of the demand letter economy. Proposals to create a registry of
demand letters would provide that vital information, allowing lawmakers and
enforcers to deal with abuses systematically.\note{see hec-testimony at 9-10}

Furthermore, many of the problems with patent demand letters stem from an
asymmetry between the sender, who has substantial legal resources and knowledge
of the patents, and the recipient, who often lacks both. Thus, requiring certain
disclosures of information in demand letters
would
put the recipients of those
letters on more even footing.\note{see hec-testimony at 11-12}
That way, discussions relating to the demand
letter can be based on the merits of the case, and not based on an imbalance in
information.

\subsection{Fair Notice to Patent Defendants}
\SectionNote

When a patent owner wishes to enforce a patent through the courts, the first
step is for the patent owner to file a paper, called a complaint, with the
court. The complaint should set forth the patent owner's basis for the case, and
put the accused infringer on notice of what is being accused.\note{see frcp at -
8/a2 (a complaint must include ``a short and plain statement of the claim
showing that the pleader is entitled to relief'')}

In every other area of the law, a complaint must provide specific details about
what the defendant did and how it injured the plaintiff. The defendant has a
right to be put on notice of the plaintiff's arguments, so that the defendant
can fairly defend the case.\note{see twombly at 559 (finding, because ``the
threat of discovery expense will push cost-conscious defendants to settle even
anemic cases,'' that ``it is only by taking care to require allegations that
reach the level suggesting conspiracy that we can hope to avoid the potentially
enormous expense of discovery in cases with no reasonably founded hope that
the discovery process will reveal relevant evidence\ldots.'')}
This is called ``notice pleading,'' and has been the
law of the United States for hundreds of years.

Due to a loophole in the rules of litigation procedure, though, patent owners
are currently exempt from this basic principle. A patent owner can simply name
defendants and a few patent numbers on the complaint, and that is enough to get
in the courtroom door.\footnote{More specifically, Form 18 of the Appendix of
Forms to the Federal Rules of Civil Procedure specifies a complaint that
essentially only identifies the patents being infringed; the Federal Circuit has
held that ``to the extent that any conflict exists between \emph{Twombly} (and
its progeny) and the Forms regarding pleadings requirements, the Forms
control.'' \sentence{k-tech at 1283; see also patentlyo-twombly}.}

Obviously, such a bare-bones complaint fails the notice requirement and denies
accused patent infringers the necessary information to mount a
defense.\note{macronix at *17 (``[Rejecting Form 18] will serve to winnow out
weak (or even baseless) claims and will protect defendants from the need to
prepare defenses for the many claims that inevitably fall by the way side in
patent cases. That also will serve to reduce the expense and burden of this kind
of litigation to both parties which, like the antitrust litigation in
\emph{Twombly}, is onerous.'')}
At a
minimum, the accused infringer deserves to know what products are being accused,
what claims of the patent are being used, and why the patent owner believes that
those products infringe the patent. Thus, a simple, straightforward fix would be
to require that information to be set forth in the complaint.

\subsection{Costs of Litigation Procedures}
\SectionNote

Patent litigation is expensive: cases can cost millions of dollars to
complete.\note{cnet-aipla (attorney fees for a fully litigated case are between
\$650,000 and \$5 million)} That cost can make it impossible for many small
businesses to defend against even an illegitimate patent: 55\% of companies sued
by patent assertion entities made \$10 million a year or less in
revenue.\note{chien-patentlyo} So the question is whether the costs of
litigation can be reduced.

Part of reducing costs is limiting the volume of patent litigation, as discussed
above.\note{see this at volume}
A patent case can be expanded to be as large as the patent owner
desires. The patent owner can augment the case easily:
\begin{itemize}
\item By accusing more products of infringement, forcing the defendant to
produce volumes of documents of product development, sales and financial
data, and arguments relating the additional products and the patents.
\item By adding more patent claims into the lawsuit, forcing the defendant to
research and interpret the additional claims, compare the accused products to
those new claims, and develop new arguments.
\item By obtaining additional patents from the Patent Office (through, for
example, patent continuation practice), thus forcing the defendant to
effectively relitigate the same case over and over again with each additional
patent.
\end{itemize}
So reducing the opportunity to expand patent litigation will help to rein in
costs.

A second potential area of cost reduction is in discovery, which is the portion
of a lawsuit in which the parties can ask each other for documents and other
information.
Discovery is an essential and valuable part of our justice system,
as it helps to ensure that both sides have access to truthful and complete
information as they prepare their cases. But discovery can be abused, and it is
often abused in patent litigation.

Abuse of the discovery process is simple: either party simply needs to demand a
large number of documents or request a large volume of information. Under the
current rules, the party responding to the requests has to pay for the costs of
those requests, so in the face of abusive discovery that party stands to outlay
enormous expenses, including time and attorney fees.

For example, a party can demand that its opponent turn over all emails relating
to a certain product and dated within the last six years. In order to comply
with that request, attorneys will have to gather and sift through all the emails
of all the employees of a company. This typically results in a database of
millions of emails. Then the attorneys will have to read through each and every
one of those millions of emails, to determine which of them include information
that is confidential or attorney-client privileged. As can easily be imagined,
this is a time-consuming, labor-intensive, expensive process.

The simple solution to patent litigation discovery abuse is limiting discovery.
Of course, care must be taken to ensure that sufficient discovery is allowed so
that the case may proceed fairly and equitably. Two specific reforms, both of
which have been proposed in bills introduced in Congress,\note{see goodlatte at section 6/a2; cornyn at section 4, S 300/a1; farenthold at section 5, S 300A}
are directed toward alleviating discovery abuses.

First, the costs
of responding to discovery can be shifted, in certain situations, to the
requesting party rather than the responding party. In most patent cases,
enormous volumes of documents are requested, but very few end up being used: one
company reported that, in one case, 10 million documents were exchanged, but
only about 2,000 were actually used at trial.\note{gao-report at 38; see also
gao-report at 37-38 (``One judge that we spoke to said that only a few of the
documents in discovery are actually used at trial---often less than one document
in 10,000\ldots.'')} Making the requesting party pay for those millions of
documents will incentivize the party to draw up narrower, more careful requests.

Second, the time
for discovery can be shifted back to later in the litigation. Specifically, a
patent case traditionally proceeds by first determining what the patent means
and then determining whether the products fit within the scope of the patent.
Determining what the patent means is often sufficient to determine the outcome
of the case, and the meaning of the patent should generally be determinable
without much discovery. Thus, the bulk of discovery could be pushed back until
after the determination of the meaning of the patent.

\subsection{Reasonable Royalty Computations}
\SectionNote
\bblabel[section \#]{reasonable-royalty}

If a court deems a patent valid, enforceable, and infringed, then the court
generally awards a ``reasonable royalty'' to the patent owner.\footnote{Where
the patent owner is an operating company, the patent owner may alternately
receive ``lost profits,'' intended to be a computation of what the patent owner
would have received absent infringement. But the reasonable royalty is always a
floor, or lower bound, on the award to a patent owner. \sentence{see 35usc at S
284}. The reasonable royalty is by far the most common award. \sentence{pwc at
11}.}
The reasonable royalty is intended to provide the patent owner with
an appropriate fraction of the revenues earned from an infringing product; in
the words of many courts, that fraction should be an amount that the patent
owner and infringer would have agreed upon in a hypothetical
negotiation.\note{see jarosz at 782-783}

Simple economics dictates that the reasonable royalty should be based on the
difference in value between the product with the patented feature, and the
product without it (or with a noninfringing alternative).\note{see, e.g.,
lemley-holdup at 1996, 1999 (defining the ``benchmark royalty rate'' as
proportional to the ``\emph{Value} per unit of the patented feature to the
downstream firm in comparison with the next best alternative technology)}
Highly useful,
inventive features would be more valuable to a product and thus would warrant a
higher royalty; small features embedded in complex, multifunction products would
warrant a lower one.\note{see chao at 122-124}

Unfortunately, the methods of computing the reasonable royalty are far from
simple. Courts often quote a legal test of fifteen different elements intended
to help determine the reasonable royalty amount.\note{see gp at 1120}
As commentators have observed,
many of these fifteen factors are duplicative or unhelpful, and the sheer number
of factors to consider often overwhelms the judges and juries tasked with
evaluating those factors.\note{see, e.g., seaman at 1703 (``Such broad,
multifactor tests have been criticized as being poorly designed and containing
duplicative or overlapping factors, which can lead to unpredictable results.'')}
Attempts to simplify this test have not helped either:
for many years a ``rule of thumb'' was to simply award 25\% of the profits of a
product, without regard to the merits of the patent at all.\note{see uniloc at
1314-1315 (``[T]his court has passively tolerated [the 25 percent rule's]
use\ldots.'')}

The complexity of assessing the reasonable royalty has led to abusive practices.
For example, because one factor considered is comparable royalty rates in the
industry, many patent owners put up websites or advertisements with inflated
royalty rates, in hopes of influencing judges and juries to adopt those numbers.
Additionally, patent owners can acquire multiple patents on the same technology,
in order to inflate the apparent value of that single technology. Some patent
owners will even sue a few weak targets in order to establish a ``market price''
for the patent.\note{see love at 1635}

Setting the royalty rate correctly is crucial to balancing the incentives that
underlie the patent system. The rate must be high enough to preserve the
incentives for inventors to disclose their inventions, but too high a rate would
stifle the marketplace, disincentivize further innovations, and undercut the
very incentives the patent system sets out to preserve.

Determining the proper
rate is a difficult task, and constructing rules to guide that determination is
even more difficult.
Shortcuts are not appropriate: even the 25\% rule was eventually rejected for
being a ``fundamentally flawed tool.''\note{uniloc at 1315 (``This court now
holds as a matter of Federal Circuit law that the 25 percent rule of thumb is a
fundamentally flawed tool for determining a baseline royalty rate in a
hypothetical negotiation.'')}
What is certain is that the existing rules must be
revised to return the reasonable royalty computation back to its basic
principles, and to avoid the sort of gamesmanship that currently goes on.

\subsection{Alternatives to Litigation}
\SectionNote

Given the high cost of patent litigation, one would hope for a low-cost
alternative to handle some of the cases where possible. And, indeed, low-cost
alternatives are provided and should be strongly supported.

The Patent Office provides several of these alternatives to litigation. Through
several procedures, the Patent Office can review issued patents to determine
whether those patents were validly issued or not, often in view of new
information such as obscure prior art not found by the patent
examiner.\footnote{\emph{Ex parte} reexamination enables a third party to ask
the Patent Office to reconsider a patent, but does not permit the third party to
participate in the proceeding. \sentence{see 35usc at S 302-307; mpep at S
2209}. Post grant review, \emph{inter partes} review, and covered business
methods review all permit the third party to participate in the proceeding; the
three procedures differ in when they may be filed, what patents they may be
filed on, and the particular questions for reconsideration that may be presented
to the Patent Office. \sentence{see 35usc at S 311-319 (inter partes review);
35usc at S 321-329 (post grant review); aia[section 18] at 329-331 (covered
business methods program); uspto-ipd}.}

These Patent Office procedures have a number of advantages. For one thing, they
are significantly lower in cost, because they deal with fewer issues and do not
require all the trappings of court litigation.\note{southard (``In light of
these figures, challenging patent validity in a district court proceeding is
likely to be a far more expensive proposition than contesting validity in the
PTO modified proceedings.'')} Additionally, the procedures are
adjudged by Patent Office employees, who generally have specialized backgrounds
in technology and patent law, as opposed to judges and juries who often lack
technical expertise and are confounded by the complexities of patent
law.\note{moore-juries at 409 (``Closer scrutiny of judge and jury
decisionmaking elucidates differences which could implicate flaws in juror
comprehension.''); moore-judges at 7 (``[F]ew district court judges are one of
ordinary skill in the technology of the invention.'')}
Furthermore, the Patent Office procedures facilitate early resolution of patent
validity questions, because many of the procedures are open to any party who
feels threatened by a patent,\note{see, e.g., 35usc at S 321/a (``[A] person who
is not the owner of a patent may file with the Office a petition to institute a
post-grant review of the patent.'')}
whereas patent litigation is only available to
those who have been formally sued or threatened with a specific
lawsuit.\footnote{The mechanism for a party to bring a district court lawsuit
against a patent owner is called a ``declaratory judgment action,'' and such an
action can only be brought where ``the facts alleged, under all the
circumstances, show that there is a substantial controversy, between parties
having adverse legal interests, of sufficient immediacy and reality to warrant
the issuance of a declaratory judgment.'' \sentence{medimmune at 127 (quoting
\clause{maryland-casualty at 273})}.}

Thus, the Patent Office procedures offer alternatives that potentially avoid the
overbearing costs of litigation. Efforts to enlarge the use of these programs
can thus further help to alleviate the potential abusive practices that arise
from those overbearing costs.

Other litigation alternatives can also be envisioned. For example, some have
proposed a small claims court for resolving smaller patent disputes on a lower
cost basis.\note{see, e.g., greenspoon} Although ideas for a patent small claims
court are generally
proposed as a way to assist small patent owners, such courts could be designed
to also assist small defendants, shielding those defendants from the abusive
discovery practices described above, for example. Similarly, some countries
provide for compulsory licensing of patents;\note{e.g., aus-patent at chapter
12, section 133} in the United States compulsory
licensing is provided for certain copyrights in music\note{e.g., 17usc at S 115
(compulsory licenses for phonorecords of nondramatic musical works)}
but not for patents.

Obviously, these litigation alternatives must be carefully considered to ensure
that they provide adequate protections to both patent owners and accused
infringers. But to solve the problem of skyrocketing patent litigation costs,
careful consideration of alternatives is warranted.

\subsection{Leveling the Playing Field}
\SectionNote

The high cost of patent litigation creates avenues for abuse, because that high
cost is borne unevenly. The owner of a patent can bring lawsuits at practically
no cost, while the defendant to that suit is guaranteed to have to pay enormous
sums no matter what the outcome. It is this imbalance in costs, to a large
degree, that has given rise to the industry of patent assertion entities,
otherwise derisively termed ``patent trolls.''\footnote{This section is
generally adapted from \sentence{cduan-forbes}.}

Patent assertion entities can structure their lawsuits so that they  have
nothing to lose when they sue others for patent infringement. Law firms will
take the cases on contingency fee arrangements, so there are no legal fees to be
paid. A successful plaintiff stands to win millions in court-awarded
royalties---a PricewaterhouseCoopers study found a median award of \$7.2 million
to non-practicing entities in 2012\note{see pwc at 7 chart 2b}---and
the unsuccessful one can just walk away.
Thus, patent assertion can be a game of all upside and no downside, of no risk
and all reward.

Contrast this with the dire situation of the defendant to a patent suit. Losing
a case means being on the hook for millions of dollars in damages. And winning
the case means being on the hook for millions of dollars in legal
fees---according to the American Intellectual Property Law Association, those
fees clock in between \$650,000 and \$5 million.\note{cnet-aipla}

The solution is to level the uneven playing field of patent law by shifting the
burden of legal fees to the loser of the lawsuit. This will place appropriate
responsibility on the patent owners, while giving victims of patent trolls a
fighting chance against weak and questionable patents. This is why fee shifting
in patent cases is part of four patent bills in Congress today.\note{shield at
section 2, S 285A/a4;
cornyn at section 5, S 285/a; goodlatte at section 3/b1, S 285/a; hatch}

But a fee shifting law alone is no better than the paper it is printed on,
because it could be easily skirted: The patent troll simply runs a business with
no assets, then when attorney fees come due, it throws up its hands, declares
bankruptcy, and walks away.

Thus, one proposed solution is to require a financial bond, ensuring that the
plaintiff actually has the assets to pay an award of attorney fees, if
necessary. Faced with an obligation to make this assurance, a patent troll with
a dubious patent will think twice before aggressively overexploiting it.

Leveling the playing field of patent litigation is an important component of
curbing abuses of the patent system. Those abuses will diminish when they
cease to be
financially viable. By undercutting the financial viability of abusive
patent assertion, one would hope that those abusers would abandon the practice
in favor of other, perhaps more beneficial, activities.

\section{Maintaining Competition in the Innovation Economy}
\SectionNote
\AuthorityTag{Maintaining Competition}

The fact that a patent is a temporary monopoly means that a patent is an
exemption from the ordinary competitive free market.\note{e.g., precisionco at
816 (``[A] patent is an exception to the general rule against monopolies and to
the right to access to a free and open market.'') (quoted in: walkerprocess at
177)}
But that fact alone does
not grant the patent owner the right to engage in all manner of anticompetitive
practices. As the Supreme Court has recognized, ``patent and antitrust policies
are both relevant in determining the `scope of the patent
monopoly.'\,''\note{actavis at 2231}

Nevertheless, less scrupulous patent owners have found methods of
abusing the patent system to engage in anticompetitive practices that go beyond
the ordinary and appropriate scope of the patent. Such anticompetitive
practices include breaking of FRAND obligations
and patent holdup, which are described below.

There are those who believe that, because the patent is a limited monopoly,
that the patent owner ought to be authorized to use patents in any manner,
regardless of the effect on competition.\note{see, e.g., actavis at 2238-2239
(Roberts, J., dissenting)} However, this misconstrues the purpose
of patents. The patent monopoly is granted not for the sake of monopoly, but as
a means to technological innovation and, ultimately, a more competitive, open
marketplace in technology. Thus, practices involving patents must be
scrutinized, as any market practices must be scrutinized, for their effects on
a competitive marketplace and consumer access to technology.\note{see actavis at
2231; see also hovenkamp at 60 (``Consumer harm does result when the
intellectual property system provides more exclusionary power than is necessary
to develop some new thing, or when it excludes without providing anything new at
all.'')}

\subsection{FRAND Obligations}
\SectionNote

This anticompetitive problem arises with regard to so-called
``standards-essential patents.'' These patents arise out of the technology
community's use of technology standards, interoperability protocols
such as WiFi or HTML, which enable different devices to communicate and operate
with each other.\note{see, e.g., ftc-doj at 33 (``Standards make networks, such
as the Internet and wireless telecommunications, more valuable by allowing
products to interoperate.'')}
Technology standards are often adopted by bodies known as
standards-setting organizations, examples of which include IEEE and
ANSI.\looseness=1

When a company develops a new technology and wishes that technology to be
incorporated into a standard, the standards-setting organization often imposes
a requirement: the company must guarantee that it will fairly license any
patents covering that technology.\note{see ftc-doj at 46-47}
This is a basic bargain: the company gets
widespread adoption of its technology, in exchange for the company agreeing to
lower license fees on patents to that technology.\note{see lemley-sep at 1137
(``FRAND commitments serve two primary goals: (1) to promote the standard by
assuring companies that they will not be blocked from bringing their products to
market\ldots, and (2) to provide reasonable rewards to those who have invested
in research and development to develop the technology used by the standard.'');
layne-farrar at 672 (``FRAND commitments are meant to address a prominent
concern in standard setting: the adoption of a technology into a major standard
could confer substantial market power, or substantially increased market power,
on its owner.'')}
This bargain is known as a
``fair, reasonable and non-discriminatory,'' or FRAND, license
obligation.\footnote{\sentence{see, e.g., layne-farrar at 672;
lemley-sep at 1160 (``Parties have spent a great deal of time litigating the
question of whether one or both sides have breached a FRAND commitment.'')}. The
term ``RAND'' is essentially interchangeable with FRAND.}\looseness=1

The anticompetitive practice arises when a holder of a FRAND-encumbered patent
breaks that obligation and begins demanding excessive royalty payments for that
patent or even attempts to block sale of the product on the market.\note{see
doj-uspto at 6 (``A patent owner's voluntary F/RAND commitments may also affect
the appropriate choice of remedy for infringement of a valid and enforceable
standards-essential patent. In some circumstances, the remedy of an injunction or exclusion order may be inconsistent with the public interest.'')}
Because the technology is already in a standard, other companies cannot
easily move away from using that technology without breaking interoperability
with other devices, thus giving the patent holder unfair leverage in the
negotiations. The contractual obligation with the standards-setting
organization might be avoided by transferring the patent to another party
unencumbered by the obligation, a practice sometimes called ``patent
privateering.''\note{see balto (``Privateering lets operating companies
evade\ldots `FRAND or other licensing commitments,' and provides a method for
`strategic outsourcing to PAEs to hinder rivals.'\,'')}

This behavior adversely affects the ability of companies to adopt new
technologies that interoperate with other products. Interoperability is a
central concern of a competitive marketplace, so the breaking of a FRAND
obligation can have serious effects on competition. Additionally, it is clear
that this anticompetitive practice does not further any interest in
compensating the inventor, for the inventor already agreed that any royalties
due under the FRAND obligation were sufficient compensation. Accordingly,
reforms directed to preventing this practice are appropriate and beneficial to
the technology marketplace.

\subsection{Patent Holdup}
\SectionNote

Many devices today include a large number of features, any one
of which could be susceptible to patenting. One study, for example, estimated
that there were 250,000 patents relevant to various aspects of
smartphones.\note{see rpx at 59}
For
such complex products and services, patents can sometimes become a ``heckler's
veto'': any one of those 250,000 patents could block access to smartphones,
particularly because a patent owner can try to obtain an injunction to block
sales of a product based on even the smallest part of that product being covered
a patent.\footnote{Several Supreme Court justices recognized the risk of
injunctions being used to demand an excessively high licensing fee.
\sentence{see ebay at 396-397 (recognizing that, for non-practicing entities,
``an injunction, and the potentially serious sanctions arising from its
violation, can be employed as a bargaining tool to charge exorbitant fees to
companies that seek to buy licenses to practice the patent.'')}.}
This problem has been described as
``patent holdup,'' because a single patent can effectively hold up production of
or access to many other unrelated technologies.\note{ftc-doj at 38 (``A holder
of IP incorporated into a
standard can exploit its position if it is costly for users of the standard to
switch to a different technology after the standard is set.''); lemley-holdup at
1992-1993 (``The threat that a patent holder will obtain an injunction that will
force the downstream producer to pull its product from the market can be very
powerful.\ldots Injunction threats often involve a strong element of
\emph{holdup} in the common circumstance in which the defendant has already
invested heavily to design, manufacture, market, and sell the product with the
allegedly infringing feature.'')}

A related problem is called ``royalty stacking.'' Where royalties are awarded
for
individual patents, a device accused of infringing multiple patents can be
subjected to multiple royalties that could theoretically reach beyond the
actual profits for the device.\note{see lemley-holdup at 1993 (``Royalty
stacking refers to situations in which a single product
potentially infringes on many patents, and thus may bear multiple royalty
burdens.'')}
For example, if a device is found to infringe 50
patents and each of those patents is adjudged to merit a 2\% royalty on revenues
from the device, then 100\% of the revenues of the device would go to patent
licenses, meaning that the device manufacturer would have to sell the device at
a loss.

Both of these problems arise out of systematic overvaluations of patents: when
courts treat patents as being more valuable than they actually are, then the
owners of those patents are able to forestall competition and innovation in the
marketplace.\note{see lemley-holdup at 2013 (``the recent surge
in patenting, especially in the information technology industry where royalty
stacking is a serious concern, these overcharges, when aggregated, can lead to a
very significant cost burden on producers.''); this at reasonable-royalty}

\section{Conclusion: The Future of Patents}

When we look back on over two hundred years of American history, we see that
patents
have played a central role in America's leadership in innovation. There is
little question that the patent system has worked well in many respects,
encouraging skilled engineers to develop new technologies and introduce those
technologies to the consumer marketplace, thereby providing society with the
benefits of advanced knowledge.

But just as we find that patents have often served as the seed of innovation,
we also find that the patent system has served as fertile ground for the
unscrupulous to take advantage of complex laws for personal benefit at the
cost of societal detriment. 

The future of patents will shape the future of innovation. In that future, we
hope that
those who invent within the patent system can work in concert with those who
work outside the patent system, to maximize the innovative capacities of both
groups. We hope that patents serve their intended purpose of disseminating
useful knowledge and are written in clear, reasonable language that gives
effective notice to all parties of what is claimed. We hope that patent holders
are compensated for their inventions, not on the basis of legal manpower and
litigious tactics, but on the basis of the merits of inventions and the value
they contribute.

And ultimately, we hope, patents will serve as a part of a competitive and fair
marketplace of technology, not geared toward transferring wealth from one party
to another but rather focused on the end goal of making technology accessible
and known to the consuming public. Such a patent system, then, would return to
its noble roots set forth in the United States Constitution, of ``promoting the
Progress of Science and the useful Arts.''

\ifshortcompile\else
\clearpage
\appendix

\twocolumn[\section{Notes}]
\markboth{}{Notes}

\theendnotes

\clearpage

\twocolumn[\section{Bibliography}]
\markboth{}{Bibliography}

\subsection*{Introduction}
\AuthorityTable{Introduction}

\subsection*{Background}
\AuthorityTable{Background}

\subsection*{Accounting for All Inventors}
\AuthorityTable{Accounting for All Inventors}

\subsection*{Clarity of Patents}
\AuthorityTable{Clarity of Patents}

\subsection*{Targeting the Right Parties}
\AuthorityTable{Targeting the Right Parties}

\subsection*{Avoiding Gamesmanship in Litigation}
\AuthorityTable{Avoiding Gamesmanship in Litigation}

\subsection*{Maintaining Competition}
\AuthorityTable{Maintaining Competition}

\fi

\clearpage
\onecolumn

\hrule height 0pt depth 0pt

\vfill

\centerline{\includegraphics[width=.5\textwidth]{pkname-large.pdf}}

\vskip2\baselineskip

\begingroup
\leftskip=.5in \rightskip=\leftskip

Public Knowledge works at the intersection of copyright, telecommunications, and
Internet law, at a time when these fields are converging. PK's experience in all
three areas puts it in an ideal position to advocate for policies that serve the
public interest.

What Public Knowledge does:

\begin{itemize}
\leftskip=.5in \rightskip=\leftskip
\item    Ensure universal access to affordable and open networks
\item    Promote creativity through balanced copyright
\item    Advance government transparency and the public’s access to knowledge
\item    Uphold and protect consumer rights
\item    Oppose policies that would slow technology, impede innovation, shrink
the public domain, or limit fair use
\item    Educate the press, the public, and policymakers using plain-language
analysis, white papers, blog posts, and videos
\item    Produce events that provide a forum for policymakers, the public,
industry, and the press to exchange ideas about our core issues
\end{itemize}

\endgroup

\vfill

\noindent Copyright \textcopyright\ 2014 Public Knowledge. This work is licensed
under the Creative Commons Attribution-ShareAlike 4.0 International License. To
view a copy of this license, visit
http://creativecommons.org/licenses/by-sa/4.0/.
Footer images:
William B. Hargrave, U.S. Pat.\ No.\ 765,361, Device for Eradicating Facial
Wrinkles, filed Nov.\ 12, 1903.
Charles A. Miller, U.S. Pat.\ No.\ 1,542,568, Radio Dial Adjuster, filed Dec.\
29, 1922.
Walter S. Olson, U.S. Pat.\ No.\ 964,803, Flying Toy Bird, filed Apr.\ 12, 1909.
Charles G. Purdy, U.S. Pat.\ No.\ 1,466,559, Exercising Device, filed Mar.\ 29,
1921.
John C. Raymond, U.S. Pat.\ No.\ 627,351, Bicycle, filed Nov.\ 4, 1897.


\end{document}
